% Options for packages loaded elsewhere
\PassOptionsToPackage{unicode}{hyperref}
\PassOptionsToPackage{hyphens}{url}
\PassOptionsToPackage{dvipsnames,svgnames,x11names}{xcolor}
\documentclass[
  english,
  10pt,
  a4paper,
]{article}
\usepackage{xcolor}
\usepackage[top=22mm,bottom=20mm,left=18mm,right=18mm]{geometry}
\usepackage{amsmath,amssymb}
\setcounter{secnumdepth}{-\maxdimen} % remove section numbering
\usepackage{iftex}
\ifPDFTeX
  \usepackage[T1]{fontenc}
  \usepackage[utf8]{inputenc}
  \usepackage{textcomp} % provide euro and other symbols
\else % if luatex or xetex
  \usepackage{unicode-math} % this also loads fontspec
  \defaultfontfeatures{Scale=MatchLowercase}
  \defaultfontfeatures[\rmfamily]{Ligatures=TeX,Scale=1}
\fi
\usepackage{lmodern}
\ifPDFTeX\else
  % xetex/luatex font selection
\fi
% Use upquote if available, for straight quotes in verbatim environments
\IfFileExists{upquote.sty}{\usepackage{upquote}}{}
\IfFileExists{microtype.sty}{% use microtype if available
  \usepackage[]{microtype}
  \UseMicrotypeSet[protrusion]{basicmath} % disable protrusion for tt fonts
}{}
\makeatletter
\@ifundefined{KOMAClassName}{% if non-KOMA class
  \IfFileExists{parskip.sty}{%
    \usepackage{parskip}
  }{% else
    \setlength{\parindent}{0pt}
    \setlength{\parskip}{6pt plus 2pt minus 1pt}}
}{% if KOMA class
  \KOMAoptions{parskip=half}}
\makeatother
\usepackage{longtable,booktabs,array}
\usepackage{caption}
\captionsetup[table]{skip=6pt}
\newcounter{none} % for unnumbered tables
\usepackage{calc} % for calculating minipage widths
% Correct order of tables after \paragraph or \subparagraph
\usepackage{etoolbox}
\makeatletter
\patchcmd\longtable{\par}{\if@noskipsec\mbox{}\fi\par}{}{}
\makeatother
\ifLuaTeX
\usepackage[bidi=basic,shorthands=off]{babel}
\else
\usepackage[bidi=default,shorthands=off]{babel}
\fi
\ifLuaTeX
  \usepackage{selnolig} % disable illegal ligatures
\fi
\setlength{\emergencystretch}{3em} % prevent overfull lines
\providecommand{\tightlist}{%
  \setlength{\itemsep}{0pt}\setlength{\parskip}{0pt}}
% ---- Çılgın Kurye document theme: same look as the CSE 396 project proposal ----
% Compile with XeLaTeX (also works on Overleaf). Fonts: TeX Gyre Termes (Times), DejaVu Sans Mono.
\usepackage{fontspec}
\setmainfont{texgyretermes}[Extension=.otf, UprightFont=*-regular, BoldFont=*-bold, ItalicFont=*-italic, BoldItalicFont=*-bolditalic]
\setsansfont{texgyreheros}[Extension=.otf, UprightFont=*-regular, BoldFont=*-bold, ItalicFont=*-italic, BoldItalicFont=*-bolditalic]
\setmonofont{DejaVuSansMono}[Extension=.ttf, UprightFont=*, BoldFont=*-Bold, ItalicFont=*-Oblique, BoldItalicFont=*-BoldOblique, Scale=0.80]
\setmathfont{texgyretermes-math.otf}

\usepackage{xcolor}
\definecolor{ckblue}{HTML}{1F3A5F}
\definecolor{cklight}{HTML}{E8EEF5}
\definecolor{ckgrey}{HTML}{F2F2F2}
\definecolor{ckline}{HTML}{555555}
\definecolor{shadecolor}{HTML}{F4F5F7}

% Headings as in the proposal: serif, bold; subsections bold italic. Numbers are part of the heading text.
\setcounter{secnumdepth}{0}
\usepackage{titlesec}
\titleformat{\section}{\rmfamily\Large\bfseries}{}{0pt}{}
\titleformat{\subsection}{\rmfamily\large\bfseries\itshape}{}{0pt}{}
\titleformat{\subsubsection}{\rmfamily\normalsize\bfseries}{}{0pt}{}
\titlespacing*{\section}{0pt}{12pt}{4pt}
\titlespacing*{\subsection}{0pt}{8pt}{3pt}

% Header and footer (the first page has only the footer, like the proposal)
\usepackage{fancyhdr,lastpage}
\pagestyle{fancy}
\fancyhf{}
\fancyhead[L]{\footnotesize\itshape CSE 396 · Group 8 · Çılgın Kurye}
\fancyhead[R]{\footnotesize\itshape\ckmodule}
\fancyfoot[C]{\footnotesize Page \thepage\ of \pageref*{LastPage}}
\renewcommand{\headrulewidth}{0.4pt}
\setlength{\headheight}{14pt}
\fancypagestyle{ckfirst}{\fancyhf{}\renewcommand{\headrulewidth}{0pt}\fancyfoot[C]{\footnotesize Page \thepage\ of \pageref*{LastPage}}}

% Inline code: breaks only after . : ( , / _ and spaces, never inside a name
\usepackage[obeyspaces,spaces]{url}
\AtBeginDocument{%
  \renewcommand{\UrlBreaks}{\do\.\do\:\do\(\do\,\do\/\do\_\do\|}%
  \renewcommand{\UrlBigBreaks}{}%
  \renewcommand{\UrlNoBreaks}{}%
  \urlstyle{tt}}

% Tables: boxed grid as in the proposal, smaller text
\usepackage{etoolbox}
\AtBeginEnvironment{longtable}{\footnotesize\renewcommand{\arraystretch}{1.2}}
\setlength{\LTpre}{6pt}
\setlength{\LTpost}{8pt}
\setlength{\arrayrulewidth}{0.4pt}

% Code blocks: long lines wrap; diagrams stay in one piece
\usepackage{framed}
\usepackage{fvextra}
\RecustomVerbatimEnvironment{verbatim}{Verbatim}{fontsize=\footnotesize, breaklines}
\makeatletter
\@ifundefined{Shaded}{}{\renewenvironment{Shaded}{\begin{snugshade}}{\end{snugshade}}}
\@ifundefined{Highlighting}{}{\RecustomVerbatimEnvironment{Highlighting}{Verbatim}{commandchars=\\\{\}, fontsize=\footnotesize, breaklines, breakanywhere}}
\makeatother

% Figures drawn with TikZ
\usepackage{tikz}
\usetikzlibrary{arrows.meta,positioning,fit,calc,backgrounds,shapes.geometric}
\tikzset{
  ck/.style={font=\footnotesize, >={Stealth[length=5pt]}, line width=0.6pt},
  box/.style={draw=ckline, fill=white, rounded corners=2pt, align=center, inner sep=4pt, minimum height=7mm},
  mod/.style={box, fill=ckgrey},
  this/.style={box, fill=cklight, draw=ckblue, line width=0.9pt},
  group/.style={draw=ckline, dashed, rounded corners=3pt, inner sep=6pt},
  lbl/.style={font=\scriptsize, align=center, inner sep=1.5pt},
  msg/.style={->, draw=ckline},
}
\usepackage{caption}
\captionsetup[figure]{font=small, labelfont=bf, skip=4pt}
\newcommand{\ckfigure}[2]{\begin{figure}[htbp]\centering #1\caption{#2}\end{figure}}

\usepackage{enumitem}
\setlist{itemsep=1pt, topsep=3pt}
\setlength{\parskip}{4pt}
\setlength{\parindent}{0pt}
\raggedbottom
% ---- end of theme ----
\newcommand{\ckmodule}{M1 Game Server and Rules}
\usepackage{bookmark}
\IfFileExists{xurl.sty}{\usepackage{xurl}}{} % add URL line breaks if available
\urlstyle{same}
% fallback for those not using the hyperref driver hyperxmp:
\makeatletter
\@ifundefined{xmpquote}{\newcommand{\xmpquote}[1]{#1}}{}
\makeatother
\hypersetup{
  pdftitle={M1 Game Server and Rules},
  pdflang={en},
  colorlinks=true,
  linkcolor={black},
  filecolor={Maroon},
  citecolor={Blue},
  urlcolor={black},
  pdfcreator={LaTeX}}

\title{M1 Game Server and Rules --- Module Documentation}
\author{}
\date{}


% ---- colours and drawings ----
\usepackage{colortbl}
\usepackage{float}
\makeatletter
\newcommand{\needspace}[1]{\begingroup\setlength{\dimen@}{#1}\vskip\z@\@plus\dimen@\penalty -100\vskip\z@\@plus -\dimen@\vskip\dimen@\penalty 9999\vskip -\dimen@\vskip\z@skip\endgroup}
\makeatother
% ---- colours ----
\definecolor{srv}{HTML}{2E6FB7}\colorlet{srvL}{srv!15}
\definecolor{cli}{HTML}{2E9E5B}\colorlet{cliL}{cli!15}
\definecolor{adm}{HTML}{E08A1E}\colorlet{admL}{adm!18}
\definecolor{brd}{HTML}{7B4FB3}\colorlet{brdL}{brd!15}
\definecolor{tcpc}{HTML}{1F5FA8}
\definecolor{udpc}{HTML}{D9531E}
\definecolor{obs}{HTML}{D7263D}
\definecolor{tA}{HTML}{00A6D6}
\definecolor{tB}{HTML}{3BB273}
\definecolor{tC}{HTML}{F2C14E}
\definecolor{tD}{HTML}{9B5DE5}

% ---- styles ----
\tikzset{
  nb/.style={rounded corners=3pt, line width=0.9pt, align=center, inner sep=4pt, minimum height=9mm, font=\footnotesize},
  sv/.style={nb, draw=srv, fill=srvL},
  cl/.style={nb, draw=cli, fill=cliL},
  ad/.style={nb, draw=adm, fill=admL},
  bd/.style={nb, draw=brd, fill=brdL},
  hwb/.style={nb, draw=black!50, fill=black!5},
  tcpa/.style={draw=tcpc, line width=1.5pt},
  udpa/.style={draw=udpc, line width=1.5pt, dash pattern=on 5pt off 2.5pt},
  inna/.style={draw=black!55, line width=1pt},
  al/.style={font=\scriptsize, fill=white, inner sep=1.2pt, rounded corners=1pt},
  grp/.style={draw=black!35, dashed, rounded corners=5pt, inner sep=7pt},
  glbl/.style={font=\scriptsize\bfseries, text=black!65},
  hl/.style={line width=2.6pt},
  st/.style={nb, rounded corners=8pt, minimum width=24mm},
  co/.style={draw=black!60, line width=0.5pt},
  cot/.style={font=\scriptsize, align=left, text=black!80},
}

% ---- system map; argument = module to highlight (M1 ... M6, or none) ----
\newcommand{\systemmap}[1]{%
\tikzset{M1/.style={},M2/.style={},M3/.style={},M4/.style={},M5/.style={},M6/.style={},none/.style={}}%
\tikzset{#1/.style={hl}}%
\begin{tikzpicture}[ck]
\node[cl, M4, minimum width=27mm] (m4) at (0,3) {\textbf{M4}\\Player phones};
\node[ad, M5, minimum width=27mm] (m5) at (0,0.4) {\textbf{M5}\\Admin panel};
\node[sv, M2, minimum width=27mm, minimum height=38mm] (m2) at (4.6,1.7) {\textbf{M2}\\Network\\and sync};
\node[sv, M1, minimum width=27mm, minimum height=38mm] (m1) at (8.4,1.7) {\textbf{M1}\\Game server\\and rules};
\node[sv, M3, minimum width=27mm, minimum height=38mm] (m3) at (12.2,1.7) {\textbf{M3}\\Physics,\\traffic\\and AI};
\node[grp, fit=(m2)(m1)(m3)] (pc) {};
\node[glbl, anchor=south west] at (pc.north west) {MAIN COMPUTER};
\node[bd, M6, minimum width=27mm, minimum height=16mm] (m6) at (4.6,-2.9) {\textbf{M6}\\ESP32 program:\\link and LED parts};
\node[hwb, minimum width=25mm] (led) at (8.4,-2.4) {RGB LED\\strips};
\node[hwb, minimum width=25mm] (sm) at (8.4,-3.6) {Sensor and\\4 motors};
\node[grp, fit=(m6)(led)(sm)] (bb) {};
\node[glbl, anchor=north east] at (bb.south east) {LED CITY BOARD};
\draw[tcpa, <->] (m4.east) -- node[al]{WebSocket} (m4.east -| m2.west);
\draw[tcpa, <->] (m5.east) -- node[al]{WebSocket} (m5.east -| m2.west);
\draw[tcpa, <->] ([xshift=-6mm]m6.north) -- node[al]{TCP} ([xshift=-6mm]m2.south);
\draw[udpa, <->] ([xshift=6mm]m6.north) -- node[al]{UDP} ([xshift=6mm]m2.south);
\draw[inna, <->] (m2) -- (m1);
\draw[inna, <->] (m1) -- (m3);
\draw[inna, ->] (m6.east |- led) -- (led);
\draw[inna, <->] (m6.east |- sm) -- (sm);
\ckleg{-1.3,-5.4}{tcpa/TCP or WebSocket: every message arrives in order, udpa/UDP: fast and the newest message wins, inna/inside one program or chip}
\end{tikzpicture}}

% ---- LED board seen from above (6 x 6 roads, 1 unit = 10 cm) ----
\newcommand{\boardbase}{%
\fill[black!88, rounded corners=3pt] (-0.5,-0.5) rectangle (5.5,5.5);
\foreach \i in {0,...,4}{\foreach \j in {0,...,4}{%
  \fill[black!60] (\i+0.13,\j+0.13) rectangle (\i+0.87,\j+0.87);
  \fill[black!22] (\i+0.24,\j+0.24) rectangle (\i+0.52,\j+0.52);
  \fill[black!30] (\i+0.56,\j+0.30) rectangle (\i+0.78,\j+0.62);
  \fill[black!38] (\i+0.28,\j+0.58) rectangle (\i+0.66,\j+0.78);}}
\foreach \i in {0,...,5}{%
  \draw[black, line width=4.5pt] (\i,-0.05) -- (\i,5.05);
  \draw[black, line width=4.5pt] (-0.05,\i) -- (5.05,\i);}
\foreach \i in {0,...,5}{\foreach \k in {0,...,30}{%
  \fill[black!55] (\i,\k/6) circle (0.025);
  \fill[black!55] (\k/6,\i) circle (0.025);}}
\foreach \x/\y/\c in {0/4/tA, 4/4/tB, 0/0/tD, 4/0/tC}{%
  \fill[\c!30!black!60] (\x+0.13,\y+0.13) rectangle (\x+0.87,\y+0.87);
  \fill[\c] (\x+0.3,\y+0.3) rectangle (\x+0.7,\y+0.7);
  \fill[white] (\x+0.4,\y+0.5) circle (0.05); \fill[white] (\x+0.6,\y+0.5) circle (0.05);}}
\newcommand{\nmark}[3]{\node[circle, fill=yellow!85!orange, draw=black, line width=0.4pt, inner sep=0.6pt, minimum size=3.6mm, font=\tiny\bfseries] at (#1,#2) {#3};}
\newcommand{\bike}[4]{%
\fill[#3!35] ($(#1,#2)+(#4:0.40)$) circle (0.05);
\fill[#3!65] ($(#1,#2)+(#4:0.21)$) circle (0.075);
\fill[#3, draw=white, line width=0.4pt] (#1,#2) circle (0.12);}
\newcommand{\dpoint}[3]{\draw[#3!50, line width=0.8pt] (#1,#2) circle (0.26); \draw[#3, line width=1.4pt] (#1,#2) circle (0.16); \fill[#3] (#1,#2) circle (0.07);}
\newcommand{\heavy}[2]{\draw[obs!60!black, line width=4pt] (#1) -- (#2); \draw[obs, line width=4pt, dash pattern=on 3pt off 5pt] (#1) -- (#2);}
\newcommand{\roadwork}[2]{\draw[obs, line width=4pt] (#1) -- (#2);}
\newcommand{\accident}[2]{\draw[obs, line width=4pt, dash pattern=on 3pt off 3pt] (#1) -- (#2);}
\newcommand{\puddle}[2]{\draw[cyan!55, line width=4pt] (#1) -- (#2); \draw[cyan!15, line width=1.6pt, dash pattern=on 1pt off 2.5pt] (#1) -- (#2);}
\newcommand{\findme}[2]{\fill[white, opacity=0.25] (#1,#2) circle (0.36); \fill[white, opacity=0.5] (#1,#2) circle (0.24); \fill[white] (#1,#2) circle (0.13);}

% ---- phone game screen (7 x 3.4 units) ----
\newcommand{\phonemock}{%
\fill[black!85, rounded corners=9pt] (-0.3,-0.3) rectangle (7.3,3.7);
\fill[cyan!20] (0,1.7) rectangle (7,3.4);
\fill[cli!14] (0,0) rectangle (7,1.7);
\filldraw[fill=srv!30, draw=srv!70!black] (0.2,1.5) rectangle (1.0,3.0);
\filldraw[fill=adm!35, draw=adm!70!black] (1.1,1.6) rectangle (1.8,2.4);
\filldraw[fill=brd!30, draw=brd!70!black] (1.9,1.65) -- (2.6,1.65) -- (2.6,2.2) -- (2.25,2.5) -- (1.9,2.2) -- cycle;
\filldraw[fill=adm!35, draw=adm!70!black] (4.5,1.65) rectangle (5.1,2.3);
\filldraw[fill=srv!30, draw=srv!70!black] (5.2,1.55) rectangle (6.0,3.1);
\filldraw[fill=cli!35, draw=cli!70!black] (6.1,1.6) rectangle (6.85,2.6);
\fill[brown!60!black] (4.17,1.6) rectangle (4.23,1.8); \fill[cli!70!black] (4.2,1.95) circle (0.18);
\fill[black!55] (1.8,0) -- (5.2,0) -- (3.75,1.7) -- (3.25,1.7) -- cycle;
\draw[white, line width=0.8pt, dash pattern=on 4pt off 4pt] (3.5,0) -- (3.5,1.7);
\filldraw[fill=black!15, draw=black!70] (3.62,1.12) rectangle (3.88,1.32);
\fill[obs] (3.05,1.22) -- (3.15,1.48) -- (3.25,1.22) -- cycle;
\fill[black!85] (3.43,0.08) rectangle (3.57,0.3);
\fill[tA] (3.3,0.28) rectangle (3.7,0.6);
\fill[brown!65] (3.34,0.6) rectangle (3.66,0.82);
\fill[tA!60!black] (3.5,0.95) circle (0.13);
\node[fill=black!70, text=white, rounded corners=2pt, font=\scriptsize\bfseries, inner sep=2pt] at (3.5,3.15) {2:41};
\foreach \c/\s/\x in {tA/120/0.35,tB/95/0.95,tC/140/1.55,tD/80/2.15}
  \node[fill=\c, text=white, font=\tiny\bfseries, rounded corners=1pt, inner sep=1.5pt, minimum width=5mm] at (\x,3.15) {\s};
\fill[white, opacity=0.9] (6.45,2.85) circle (0.38); \draw[black!60] (6.45,2.85) circle (0.38);
\draw[-{Stealth[length=5pt]}, line width=1.8pt, obs] (6.45,2.85) ++(230:0.2) -- ++(50:0.5);
\fill[white, opacity=0.85] (0.55,0.5) circle (0.32); \fill[black!70] (0.43,0.5) -- (0.65,0.63) -- (0.65,0.37) -- cycle;
\fill[white, opacity=0.85] (1.4,0.5) circle (0.32); \fill[black!70] (1.52,0.5) -- (1.3,0.63) -- (1.3,0.37) -- cycle;
\node[fill=obs!80, text=white, rounded corners=3pt, font=\tiny\bfseries, minimum width=8mm, minimum height=6mm] at (5.6,0.5) {Brake};
\node[fill=cli, text=white, rounded corners=3pt, font=\tiny\bfseries, minimum width=8mm, minimum height=8mm] at (6.5,0.6) {Gas};
\node[fill=white, text=black, rounded corners=3pt, font=\tiny\bfseries, minimum width=9mm, minimum height=5mm] at (0.9,1.45) {Find me};}

% ---- side cut of the board (14 units wide) ----
\newcommand{\boardside}{%
\fill[brown!20] (-0.6,-0.35) rectangle (14.6,0);
\foreach \x/\h/\c in {1.0/0.7/black!18, 2.2/1.1/black!22, 4.6/1.9/black!25, 5.8/2.6/black!28, 7.4/2.2/black!25, 8.6/1.6/black!22, 10.9/1.0/black!18}
  \filldraw[fill=\c, draw=black!45] (\x,1.5) rectangle (\x+0.9,1.5+\h);
\filldraw[fill=tB!35, draw=tB!70!black] (12.4,1.5) rectangle (13.4,2.3);
\fill[black!70, rounded corners=1pt] (0.5,0) rectangle (1.1,0.9);
\fill[black!70, rounded corners=1pt] (12.9,0) rectangle (13.5,0.9);
\filldraw[fill=brown!45, draw=brown!70!black] (0,0.9) rectangle (14,1.3);
\fill[black!80] (0,1.3) rectangle (14,1.38);
\foreach \x in {0.3,0.9,...,13.8} \fill[black!35] (\x,1.42) circle (0.06);
\foreach \x/\c in {3.3/tA, 3.9/tA, 6.9/obs, 7.5/obs, 10.5/tC} \fill[\c] (\x,1.42) circle (0.07);
\fill[white, opacity=0.55] (0,1.38) rectangle (14,1.52);
\filldraw[fill=brd!40, draw=brd] (2.3,0.35) rectangle (4.3,0.85);
\node[font=\tiny] at (3.3,0.6) {ESP32};
\filldraw[fill=brd, draw=brd!60!black] (6.75,0.62) rectangle (7.25,0.88);
\filldraw[fill=adm!60, draw=adm!70!black] (1.8,0.6) circle (0.24);
\filldraw[fill=adm!60, draw=adm!70!black] (12.2,0.6) circle (0.24);
\filldraw[fill=black!15, draw=black!50] (9.3,0.35) rectangle (10.9,0.85);
\node[font=\tiny] at (10.1,0.6) {Power};}

% ---- layout helpers ----
\newenvironment{ckfig}{\begin{figure}[!htb]\centering}{\end{figure}}
\renewcommand{\topfraction}{0.9}\renewcommand{\bottomfraction}{0.8}\renewcommand{\textfraction}{0.07}\renewcommand{\floatpagefraction}{0.85}
\setcounter{topnumber}{3}\setcounter{bottomnumber}{2}\setcounter{totalnumber}{4}
\newcommand{\msub}[1]{\par\needspace{5\baselineskip}\smallskip\noindent{\bfseries #1:}\par\nobreak\vspace{-2pt}}
\tikzset{msga/.style={draw=cli!70!black, line width=0.8pt}, cmda/.style={draw=obs, line width=1pt}, loopa/.style={draw=srv, line width=1pt}, hwa/.style={draw=black!60, line width=1.2pt, dash pattern=on 1.5pt off 1.5pt}, wanta/.style={draw=cli, line width=1.6pt}, reala/.style={draw=obs, line width=1.6pt}, siga/.style={draw=srv, line width=0.8pt}, thra/.style={draw=obs, line width=0.9pt, dash pattern=on 4pt off 2pt}, evta/.style={draw=cli, line width=1pt}}
\tikzset{iol/.style={font=\scriptsize, fill=white, inner sep=1.2pt, rounded corners=1pt, align=center}}
% legend row: start point, list of style/text
\newcommand{\ckleg}[2]{%
\coordinate (lgp) at (#1);
\foreach \ls/\lt in {#2}{%
  \draw[\ls,->] (lgp) -- ++(0.8,0) node[right, font=\scriptsize, text=black] (lgt) {\lt};
  \coordinate (lgp) at ([xshift=5mm]lgt.east);}}
% purpose table: colour (unused), runs on, technology, requirements
\newcommand{\ckinfo}[4]{\par\smallskip\noindent{\renewcommand{\arraystretch}{1.25}%
\begin{tabular}{|L{0.30}|L{0.50}|L{0.20}|}\hline
\textbf{Runs on} & \textbf{Technology} & \textbf{Requirements} \\ \hline
#2 & #3 & #4 \\ \hline
\end{tabular}}\par\smallskip}
% module data map: n in, n out, centre text, centre style, inputs, outputs, legend list
% inputs and outputs: {name}/style/{data}/arrow style
\newcommand{\iomap}[7]{%
\begin{tikzpicture}[ck]
\pgfmathsetmacro{\nmax}{max(#1,#2)}
\node[#4, line width=2pt, minimum width=30mm, minimum height=\nmax*10.5mm+2mm, font=\normalsize] (c) at (0,0) {#3};
\foreach \nm/\st/\dd/\ar [count=\ii] in {#5}{%
  \pgfmathsetmacro{\yy}{((#1+1)/2-\ii)*1.05}
  \node[\st, minimum width=25mm, minimum height=8.5mm, font=\scriptsize] (in\ii) at (-6.6,\yy) {\nm};
  \draw[\ar,->] (in\ii.east) -- node[iol]{\dd} (c.west |- in\ii);}
\foreach \nm/\st/\dd/\ar [count=\ii] in {#6}{%
  \pgfmathsetmacro{\yy}{((#2+1)/2-\ii)*1.05}
  \node[\st, minimum width=25mm, minimum height=8.5mm, font=\scriptsize] (out\ii) at (6.6,\yy) {\nm};
  \draw[\ar,->] (c.east |- out\ii) -- node[iol]{\dd} (out\ii.west);}
\node[font=\scriptsize\bfseries, text=black!60] at (-6.6,\nmax*0.525+0.4) {Inputs come from};
\node[font=\scriptsize\bfseries, text=black!60] at (6.6,\nmax*0.525+0.4) {Outputs go to};
\ckleg{-7.6,-\nmax*0.525-0.5}{#7}
\end{tikzpicture}}

% keep an intro sentence and its table together on one page
\newcommand{\ckkeep}[1]{\par\noindent\begin{minipage}{\linewidth}\setlength{\parskip}{4pt}#1\end{minipage}\par\smallskip}
\widowpenalty=10000 \clubpenalty=10000

\begin{document}

\thispagestyle{ckfirst}
\begin{center}
{\bfseries GEBZE TECHNICAL UNIVERSITY\par}
{Computer Engineering Department\par}
{CSE 396 Module Documentation\par}
\vspace{8pt}
{\LARGE\bfseries Çılgın Kurye\par}
\vspace{3pt}
{\large\itshape Module M1: Game Server and Rules\par}
\vspace{6pt}
\textbf{Company:} Board\&Byte\hspace{1.4em}\textbf{Group:} Group-8\hspace{1.4em}\textbf{Date:} 7 November 2026
\end{center}
\begin{center}\footnotesize\renewcommand{\arraystretch}{1.15}
\begin{tabular}{|p{0.27\linewidth}|p{0.15\linewidth}|p{0.27\linewidth}|p{0.15\linewidth}|}\hline
\textbf{Name} & \textbf{Student No.} & \textbf{Name} & \textbf{Student No.} \\ \hline
Tuğberk Akbulut & 230104004070 & Abdullatif Yılmaz & 230104004031 \\ \hline
Berk Andaçgil & 230104004066 & Selim Barut & 220104004046 \\ \hline
Umutcan Ocak & 220104004903 &  &  \\ \hline
\end{tabular}
\end{center}
\vspace{2pt}



% ---- table and list helpers ----
\newcolumntype{L}[1]{>{\raggedright\arraybackslash}p{\dimexpr#1\linewidth-2\tabcolsep-\arrayrulewidth\relax}}
\newcolumntype{C}[1]{>{\centering\arraybackslash}p{\dimexpr#1\linewidth-2\tabcolsep-\arrayrulewidth\relax}}
\newcommand{\ckdone}{\item[$\boxtimes$]}
\newcommand{\ckopen}{\item[$\square$]}
\newcommand{\mdot}{$\bullet$}
\newcounter{msec}
\newcounter{chap}
% ---- end of helpers ----
\newcommand{\msection}[1]{\needspace{8\baselineskip}\stepcounter{msec}\section{\arabic{msec}. #1}}

\msection{Purpose}

M1 is the brain of the game. It runs on the main computer and it is the only part that can change
the game. It keeps the game clock and the one true game state. It knows every team, store, package,
delivery point, score, fall, obstacle and earthquake. Phones only send their buttons. M1 decides
what happens, and then everybody sees the same result.

\ckinfo{srv}{Main computer, in one server program together with M2 and M3}{TypeScript on Node.js, a fixed step game loop with 30 ticks per second, seeded random numbers, JSON files for the city plan and the game log}{R1 R2 R4 R5 R6 R7 R9 R11}

M1 has no network code. All its inputs and outputs are data in the memory of the server program. The
figure shows which module gives M1 each kind of data and which module gets each result.

\begin{ckfig}
\iomap{6}{6}{\textbf{M1}\\Game server\\and rules}{sv}
{{M2, from phones}/cl/{controls, joins,\\``Find me''}/inna, {M2, from admin}/ad/{commands,\\parameters}/inna, {M2, from board}/bd/{earthquake\\events}/inna, {M3 Traffic}/sv/{NPC cars, bot controls,\\travel times}/inna, {M3 Physics}/sv/{poses, collisions,\\falls}/inna, {City plan file}/hwb/{roads, stores,\\delivery points}/inna}
{{M2, to phones}/cl/{game state\\30 per second}/inna, {M2, to admin}/ad/{state, metrics,\\replies}/inna, {M2, to board}/bd/{board state,\\vibrate}/inna, {M3 Physics}/sv/{controls, obstacles,\\earthquake}/inna, {M3 Traffic}/sv/{players, obstacles,\\parameters}/inna, {Disk}/hwb/{game log}/inna}
{inna/data in the memory of the server program}
\end{ckfig}

\msection{Requirements covered}

\ckkeep{%
M1 takes part in eight of the eleven system requirements. The last column says what will work when
M1 is finished.

\begin{tabular}{|L{0.06}|L{0.15}|L{0.40}|L{0.39}|}
\hline \textbf{ID} & \textbf{Requirement} & \textbf{What M1 does for it} & \textbf{When M1 is done} \\ \hline 
R1 & Multiplayer game & Lobby with 4 teams and at most 2 players in a team. Bots fill the empty places. A game is at most 5 minutes. & A full game runs from the lobby to the winner, with humans and bots together \\ \hline
R2 & Delivery cycle & Automatic pickup at the store. Picks a delivery point by travel time, with a near or far chance per player & Every player always has a package or a point to go to, and points are fair \\ \hline
R4 & Collisions and falls & Gives a penalty time after every fall that M3 reports & A fallen player waits and then rides on with the package \\ \hline
R5 & No map on the phone & Puts the compass target and the ``Find me'' waiting time into the game state & The phone can show the arrow and the button without a map \\ \hline
R6 & Obstacles & Keeps the one obstacle list that M3, M4 and M6 use. Adds random obstacles during the game and removes each obstacle when its time ends & A new or ended obstacle, automatic or from the admin, is seen everywhere in the same tick \\ \hline
R7 & Earthquake & One path for a hand hit, the motors and the admin button. Road work and falls inside the radius & A hit on the board changes the game for everybody \\ \hline
R9 & Live admin panel & Checks admin commands and sends each parameter to the module that uses it & Every parameter can be changed during the game \\ \hline
R11 & Engineering requirements & Only M1 changes the game state. Every game can be played again from the log & The same log always gives the same game \\ \hline
\end{tabular}}

\msection{Inputs and outputs}

\ckkeep{%
The table gives the content of each input and output. Value ranges are shown only for the main
control values.

\begin{tabular}{|L{0.11}|L{0.17}|L{0.18}|L{0.54}|}
\hline \textbf{Direction} & \textbf{Data} & \textbf{From / to} & \textbf{Content} \\ \hline 
In & Controls & M2, from phones & Throttle and brake from 0 to 1, steer from −1 to 1, a control number. ``Find me'' comes as an event \\ \hline
In & Join & M2, from phones & Player key, name, team. M2 also reports when a player is lost or comes back \\ \hline
In & Admin commands & M2, from admin & Parameter changes, start, pause, resume, end, reset, new game code, obstacles, earthquake \\ \hline
In & Board events & M2, from board & Source (hand or motors), strength from 0 to 10, zone, confidence \\ \hline
In & Traffic & M3 & NPC cars, bot controls, travel time table \\ \hline
In & Physics & M3 & New motorcycle poses, collision and fall events \\ \hline
Out & Game state & M2, to phones and admin & Every tick: phase, time, scores, players, NPC cars, obstacles, earthquake. For each player also the compass target, the ``Find me'' waiting time and the number of the last control that M1 used \\ \hline
Out & Board state & M2, to board & Phase, players, delivery points, obstacles and earthquake in board coordinates \\ \hline
Out & Metrics & M2, to admin & Twice per second: the values of all admin graphs \\ \hline
Out & Game log & Disk & The seed and, for every tick, the controls, admin commands and board events \\ \hline
\end{tabular}}

\msection{Main components}

\ckkeep{%
M1 is made of seven parts. Each part takes some data and gives a result to the next part or to
another module.

\begin{tabular}{|L{0.21}|L{0.39}|L{0.40}|}
\hline \textbf{Component} & \textbf{Takes} & \textbf{Gives} \\ \hline 
Game loop & The clock & Runs the steps of a tick in a fixed order, 30 times per second \\ \hline
Match and lobby & Joins, lost players, admin start and pause & Game phase, teams, bots in empty places \\ \hline
Delivery and scoring & Poses, travel times & Packages, delivery points, automatic deliveries, points \\ \hline
Penalties & Fall events from M3 & Penalty time, then a stand up request to M3 \\ \hline
Obstacles and earthquake & Obstacle parameters, admin obstacles, board events & The obstacle list with end times, earthquake centre and radius \\ \hline
Admin and parameters & Admin commands & Checked values sent to the owner module, ok or error replies \\ \hline
State builder and recorder & The world after the tick & Frozen game state, board state, metrics with a short history, the game log \\ \hline
\end{tabular}}

\msection{How it works}

\msub{What a tick is}
The game world does not change all the time. It changes in small steps, like the pictures of a film.
One step is called a tick, and the server makes 30 ticks in one second. Between two ticks, M2 only
collects the messages from the phones, the admin and the board. In a tick, M1 uses all of them
together, moves the world one step forward and makes a new state. Then M2 sends this state to
everybody. So every player sees the same world after every tick.

\begin{ckfig}
\begin{tikzpicture}[ck]
\draw[black!45, ->, line width=0.8pt] (-0.3,0) -- (14.6,0) node[right, font=\scriptsize]{time};
\foreach \x/\l in {1.2/1, 5.8/2, 10.4/3}{%
  \node[sv, minimum width=17mm, minimum height=8mm, font=\scriptsize] (t\l) at (\x,0) {Tick \l};
  \draw[inna, ->] (t\l.north) -- ++(0,0.8) node[above, font=\scriptsize, align=center]{new state\\to M2};}
\foreach \qa/\qb/\n in {2.3/4.3/2, 6.9/8.9/3}{%
  \node[nb, draw=cli, fill=cliL, minimum width=20mm, minimum height=7mm, font=\scriptsize] (q\n) at ({(\qa+\qb)/2},-1.5) {};
  \foreach \k in {0,1,2} \fill[cli] ({\qa+0.45+\k*0.55},-1.5) circle (0.09);
  \node[font=\scriptsize, text=cli!70!black] at ({(\qa+\qb)/2},-2.15) {messages wait in the queue};
  \draw[msga, ->] (q\n.east) -| (t\n.south);}
\draw[srv, <->, line width=0.9pt] (1.2,2.25) -- node[iol]{one tick: one thirtieth of a second} (5.8,2.25);
\ckleg{0,-2.85}{inna/new state after each tick, msga/{messages from phones, admin and board}}
\end{tikzpicture}
\end{ckfig}

\msub{The steps of one tick}
A tick always runs the same eight steps in the same order. First M1 takes the inputs that M2 has
collected. Then it applies admin commands and board events, updates the game phase, asks M3 for
traffic and bot controls, and then asks it to move all motorcycles. After that the game rules run. At the
end M1 builds a frozen copy of the state for M2 and writes one line to the game log. Because the order
never changes and all random numbers come from one seed, the same inputs always give the same game.

\begin{ckfig}
\begin{tikzpicture}[ck, node distance=3.5mm]
\node[cl, minimum width=17mm] (a) {Inputs};
\node[ad, minimum width=17mm, right=of a] (b) {Admin and\\board events};
\node[sv, minimum width=17mm, right=of b] (c) {Game\\phase};
\node[sv, minimum width=17mm, right=of c, fill=tB!20, draw=tB!80!black] (d) {Traffic\\M3};
\node[sv, minimum width=17mm, right=of d, fill=tD!18, draw=tD!80!black] (e) {Physics\\M3};
\node[sv, minimum width=17mm, right=of e, fill=obs!12, draw=obs] (f) {Rules};
\node[sv, minimum width=17mm, right=of f] (g) {Frozen\\state};
\node[hwb, minimum width=17mm, right=of g] (h) {Log};
\foreach \u/\v in {a/b,b/c,c/d,d/e,e/f,f/g,g/h} \draw[inna,->] (\u) -- (\v);
\draw[inna,->, srv] (h.south) -- ++(0,-0.6) -| node[al, pos=0.25]{next tick} (a.south);
\ckleg{0,-1.6}{inna/order of the steps, loopa/start of the next tick}
\end{tikzpicture}
\end{ckfig}

\msub{Controls inside a tick}
A phone sends its controls on every change and at least 10 times per second, but there are 30 ticks
per second. So in some ticks a new control comes and in some ticks nothing comes. If nothing comes,
M1 keeps the last control of that player. If more than one control comes, M1 uses the newest one for
throttle, brake and steer. ``Find me'' is different. It is an event, not a value. If any message in
the tick has a ``Find me'' press, the press counts, even when a newer message has no press. So a
button press is never lost.

\msub{Game flow}
A game has five phases. In the lobby players join and choose a team. When the admin presses start, a
3 second countdown begins and nobody can move. Then the game runs for at most 5 minutes. The admin can
end the game early. The admin can also pause and resume at any time. During a pause everything stops:
the game time, the penalty times, the waiting time for lost players and the NPC cars. A pause during
the countdown starts the countdown again after the resume. When the time ends, the team with the most
points wins. If two teams have the same points, the game ends in a draw between them.

\begin{ckfig}
\begin{tikzpicture}[ck]
\node[st, draw=cli, fill=cliL] (l) at (0,0) {Lobby};
\node[st, draw=adm, fill=admL] (c) at (3.6,0) {Countdown};
\node[st, draw=srv, fill=srvL, line width=1.4pt] (r) at (7.2,0) {Running};
\node[st, draw=brd, fill=brdL] (f) at (11.2,0) {Finished};
\node[st, draw=black!50, fill=black!6] (p) at (5.4,-2.0) {Paused};
\draw[inna,->] (l) -- node[al]{start} (c);
\draw[inna,->] (c) -- node[al]{3 s} (r);
\draw[inna,->] ([yshift=1.5mm]r.east) -- node[al, above=1pt]{time over} ([yshift=1.5mm]f.west);
\draw[inna,->, obs] ([yshift=-1.5mm]r.east) -- node[al, below=1pt]{admin end} ([yshift=-1.5mm]f.west);
\draw[inna,->] ([xshift=-6mm]r.south) |- node[al, pos=0.3]{pause} ([yshift=1.5mm]p.east);
\draw[inna,->] ([yshift=-1.5mm]p.east) -| node[al, pos=0.7]{resume} ([xshift=6mm]r.south);
\draw[inna,->] (c.south) |- node[al, pos=0.3]{pause} (p.west);
\draw[inna,->] (f.north) -- ++(0,0.6) -| node[al, pos=0.25]{new game} (l.north);
\ckleg{0.2,-3.0}{inna/phase change, cmda/admin command}
\end{tikzpicture}
\end{ckfig}

\msub{Players who leave and come back}
M2 sees when the link of a phone is lost and tells M1. M1 keeps the place of this player for 30
seconds. If the player does not come back in this time, a bot takes the motorcycle and plays on. If
the player comes back later with the same player key, the bot gives the motorcycle back. The admin
sets how many bots fill the empty places at the start.

\msub{Deliveries and points}
A player gets a package automatically when the motorcycle is inside the area of the own team's store.
M1 then picks a delivery point for this player. M3 gives a table with the travel time from the store
to every possible point. The times count the obstacles that are on the roads at that moment. The
admin sets for each player whether the new delivery point appears near the store or far from it, and
M1 picks the point with this setting. Two players never get the same point at
the same time. A far point gives more points. Every delivery gets a base value, and the longer the
expected travel time is, the bigger the bonus on top of it. The admin sets both the base value and
the bonus. The points are fixed when the point is chosen, so a new obstacle
later does not change them.

There is no delivery button. When the player comes to the delivery point and stops inside its area,
M1 delivers the package automatically and adds the points. The size of the area is an admin
parameter. A fall does not lose the package. The player only waits for the penalty time.

\begin{ckfig}
\begin{tikzpicture}[x=1.1cm, y=1.1cm]
\boardbase
\draw[tA, line width=2.4pt, rounded corners=3pt, opacity=0.9] (1,4) -- (1,3) -- (2,3) -- (2,2) -- (2.5,2);
\dpoint{2.5}{2}{tA}
\draw[tA!45, line width=1pt, dashed] (4,1.5) circle (0.18);
\bike{1}{3.5}{tA}{90}
\bike{3.5}{4}{tB}{0}
\bike{4}{2.5}{tC}{-90}
\bike{1.5}{1}{tD}{180}
\roadwork{3,2}{3,3}
\nmark{0.5}{5.2}{1}\nmark{0.55}{3.5}{2}\nmark{1.5}{2.7}{3}\nmark{2.5}{2.5}{4}\nmark{4.45}{1.5}{5}\nmark{3.35}{2.5}{6}
\node[anchor=west, font=\small, align=left] at (6.0,2.5) {\begin{tabular}{@{}c@{\ \ }l@{}}
\textbf{1} & Team store in the team colour\\[3pt]
\textbf{2} & Player with a package,\\ & a dot with a tail\\[3pt]
\textbf{3} & Fastest route found by M3\\[3pt]
\textbf{4} & Near delivery point,\\ & blinks in the player's colour\\[3pt]
\textbf{5} & A far point, worth more points\\[3pt]
\textbf{6} & Road work, so the route\\ & goes around\end{tabular}};
\end{tikzpicture}
\end{ckfig}

\msub{Earthquakes}
An earthquake can come from a hand hit on the board, from the vibration motors or from the admin
button. All of them take the same path in M1. The board sends a strength from 0 to 10 and one of
9 zones. M1 uses the centre of that zone as the centre of the earthquake. A stronger hit makes a
bigger earthquake. Motorcycles near the centre fall, and the roads near the centre close with road
work for a limited time.

\msub{Obstacles}
M1 keeps one list of obstacles. Each obstacle has a type, a road segment, a lane, a position and an end
time. During the game M1 adds new obstacles by itself at random places. The admin sets for each of the
four types how often it appears and how long it stays. A new obstacle never appears right next to a
player or right in front of one, so nobody falls without a chance to react. Obstacles from the admin
and the road work from an earthquake also have a duration. When the time of an obstacle ends, M1
removes it from the list. While the game is paused, no new obstacle appears and the times stop. When
the list changes, M3 computes the routes again, and M4 and M6 show the change.

\msub{Admin parameters}
Every parameter has a name, a group, a minimum, a maximum and an owner module. When the admin changes a value, M1 checks the limits first.
Then it saves the value, sends it to the owner (for example grip to M3) and sends ``ok'' back to the
panel. A wrong value gets an error and the old value stays.

\msection{Testing}

M1 is tested alone, without the network and without the real M2 and M3. Small fake parts take
their place inside the test and give fixed test data: scripted controls and joins, fixed poses and fall
events, a fixed travel time table and recorded admin commands and board events. With this data we
check every rule: the game phases with pause and end, the lobby and the bots, a lost player who comes
back, the delivery cycle and the fair choice of points, two ``Find me'' presses in one tick, every
admin command with good and wrong values, the earthquake radius, and the automatic obstacles with
their rates, durations, distance to the players and the stop during a pause. Because all random numbers come
from one seed, the same test data must always give exactly the same game state.

\end{document}
